\originalchapter{8.333复习包2：相空间、等离子体与输运}
\osource{8.333 Fall 2013页面公开的Fall 2007 Test2 Review及Solutions，7题。对Coulomb势、截面角度和Lorentz气Wick系数作独立符号与量纲复核，修正处明确写出。}
\begin{exercise}[原题1(a)--(d)]\label{r:333-2-1}
一维自由气粒子初始$\rho(q,p,0)=\delta(q)f(p)$，$f=\exp[-p^2/(2m\kb T)]/\sqrt{2\pi m\kb T}$。(a)Liouville解和相图；(b)两方差；(c)在$q=\pm Q$加硬壁，长时细密度；(d)有限相空间分辨率的粗粒化密度及其平稳性。
\end{exercise}
\begin{solution}
自由特征$q(t)=q_0+pt/m$、$p$常量，$\rho=\delta(q-pt/m)f(p)$，支持是一条随$t$倾斜的直线而非扩散成二维Gaussian。$\avg{q^2}=\kb Tt^2/m$、$\avg{p^2}=m\kb T$。硬壁反射$\pm p$，将直线按周期$4Q$折回盒子，得到细密度越来越细的折线，仍保存每一相点信息和$|p|$。

对固定小$q,p$单元平均，足够宽的速度小区间在长时覆盖许多反射周期，弱/粗粒化极限$\widetilde\rho=f(p)/(2Q)$（初始$f$偶）。它在盒内$q$均匀且满足反射对称，故Liouville平稳；精细密度本身不点态收敛。初始理想$\delta$的微分熵需先有限宽正则化，不能以无穷负熵作有限数值热力学熵。
\end{solution}
\begin{exercise}[原题2(a)--(d)]\label{r:333-2-2}
相空间规范化密度$S=-\int\rho\log\rho$。(a)Hamilton--Liouville演化$\dot S$；(b)固定$\avg H=E$的最大熵；(c)其平稳性；(d)如何与观察到趋平衡的熵增相容。
\end{exercise}
\begin{solution}
$\dot\rho=-\{\rho,H\}$，Hamilton流无散度且边界通量0，$\dot S=\int(1+\log\rho)\{\rho,H\}=\int\{\rho\log\rho,H\}=0$。对规范化和能量加乘子，变分$-(1+\log\rho)-a-\beta H=0$，故$\rho_{\max}=\e^{-\beta H}/Z$，$\beta$由$E$定；熵对$\rho$凹，确保最大。$\{f(H),H\}=0$，canonical平稳，但不意味着任何初态自动变成它。

粗粒化丢弃细结构，由凹性其熵不小于细密度熵；局部混合、与外界交换、或Boltzmann的分子混沌假设可使可观察熵增加。完全隔离精密可逆Hamilton演化并没有违反(a)，而coarse熵也未由Liouville单独保证每一时刻严格单调。
\end{solution}
\begin{exercise}[原题3(a)--(c)]\label{r:333-2-3}
Vlasov适用于相互作用范围$\lambda$内粒子数$n\lambda^3\gg1$，忽略相对弱的短时两体碰撞项。(a)$\rho_N=\prod_i\rho_1(x_i)$的$f_s$与规范化；(b)从BBGKY得$\partial_tf_1+(p/m)\cdot\nabla_qf_1-\nabla_qU_{\rm eff}\cdot\nabla_pf_1=0$，$U_{\rm eff}=U+\int V(q-q')f_1(x')\dd x'$；(c)均匀盒内任意动量$g(p)$为何平稳而无热化。
\end{exercise}
\begin{solution}
取$\int\rho_1=1$、约化密度按计数规范化，$f_s=N!/(N-s)!\prod_{i=1}^s\rho_1(x_i)$，$\int f_s=N!/(N-s)!$。$f_2=(1-1/N)f_1f_1$，代BBGKY的一体式得到平均场势$(1-1/N)\int Vf_1$，在大$N$忽略自作用的相对$1/N$差便是题式；各$s$乘积按同一一体流因子分离。这个无相关闭合不是任意有限$N$两体系统的精确演化。

均匀$f_1=N g(p)/V$（$\int g=1$）的密度势为空间常数，流项全0，所以任意$g$平稳。若原题把归一化单粒子密度写$g/V$，与计数$f_1$差一个$N$。没有碰撞核就没有要求$g$成为Maxwell的H定理；可有集体相混合，却不等于完整碰撞热平衡。
\end{solution}
\begin{exercise}[原题4(a)--(f)]\label{r:333-2-4}
中性等离子体各$N$个$\pm e$、质量$m_\pm$、密度$n_0$。(a)两组Vlasov及平均Coulomb势；(b)$f_\pm=g_\pm(p)n_\pm(q)$的Poisson式；(c)Boltzmann密度与Debye长；(d)屏蔽Green解；(e)Vlasov自洽条件及平均间距；(f)$\tau\sim\lambda/c_{\rm th}$为何不随$T$变、关联何物理频率。原题$\nabla^2C=+4\pi\delta$的势约定与正$1/r$Green式不相容，以下先统一标准物理势$\phi$。
\end{exercise}
\begin{solution}
取$C(r)=1/r$、$\nabla^2C=-4\pi\delta$、粒子势能$q_s\phi$，则$[\partial_t+(p/m_s)\cdot\nabla-q_s\nabla\phi\cdot\nabla_p]f_s=0$，$\phi=\phi_{ext}+e\int C(f_+-f_-)\dd x'$。规范化$\int g_s\dd p=1$给$\nabla^2\phi=-4\pi[\rho_{ext}+e(n_+-n_-)]$。

热平衡$n_\pm=n_0\e^{\mp\beta e\phi}$，弱势线性$e(n_+-n_-)\simeq-2\beta n_0e^2\phi$，于是$(\nabla^2-\lambda^{-2})\phi=-4\pi\rho_{ext}$、$\lambda^2=\kb T/(8\pi n_0e^2)$。解$\phi(r)=\int\e^{-|r-r'|/\lambda}\rho_{ext}(r')/|r-r'|\dd^3r'$，因为$(\nabla^2-\lambda^{-2})(\e^{-r/\lambda}/r)=-4\pi\delta(r)$；仅在$r\ne0$验证微分式会漏掉决定符号的点源。若使用原题$\Phi=-\phi$、$C=-1/r$则Vlasov正负流及$n_\pm=n_0\e^{\pm\beta e\Phi}$转换一致，但Green要取负号。

$n_0\lambda^3\gg1$等价$\lambda\gg\ell=n_0^{-1/3}$，亦即Coulomb能$e^2/\ell\ll\kb T$（常数阶略去），很多粒子共同形成弱平均场。$c_{\rm th}\sim\sqrt{\kb T/m}$，$\lambda/c_{\rm th}\sim\sqrt{m/(n_0e^2)}\sim\omega_p^{-1}$，温度抵消。这是集体响应/振荡时间尺度，碰撞热化或Landau阻尼时间未必同一个常数。
\end{solution}
\begin{exercise}[原题5(a)--(i)]\label{r:333-2-5}
二维经典电子均匀磁场，单体$H=(p-q_eA)^2/(2m)\pm\mu_B|B|$、$A=B\times q/2$，分自旋上、下，守自旋碰撞截面$\sigma_{\uparrow\uparrow}=\sigma_{\downarrow\downarrow}=\sigma,\sigma_{\uparrow\downarrow}=\sigma_\times$。(a)轨道；(b)两Boltzmann式；(c)总H单调；(d)碰撞零项条件；(e)流项的守恒量函数；(f)角动量碰撞内外；(g)圆对称约束的最一般平衡分布；(h)磁/非磁杂质影响；(i)自旋与角动量是否给额外流体方程。
\end{exercise}
\begin{solution}
机械动量$\pi=p-q_eA$，$\dot q=\pi/m$、$\dot\pi=q_e(\pi/m)\times B$，横向回旋频率$|q_eB|/m$，均匀Zeeman常数不改轨道。每组的流算子$L_s$加$C_{ss}$与$C_{s\bar s}$；具体$C_{ss'}=\int\dd\pi_2\dd\Omega\,|v-v_2|\sigma_{ss'}[f_s'f_{s',2}'-f_sf_{s',2}]$，上、下互换得第二式，局域同位置碰撞。

乘$1+\log f_s$求和，把反向散射和两入射粒子交换平均，$\dot H=-\tfrac14\sum_{ss'}\int W(A-B)\log(A/B)\le0$。等号要求$\log f_s$为碰撞不变量线性组合，即两组各自数乘子、共享总动量与能量系数。若还要求平稳流，乘子组合需是单体Hamilton守恒量；任何$f(H_s,L_z)$的Poisson流可为0，但只有合适指数线性组合同时满足碰撞条件。

圆对称时\textbf{总canonical}$L_z=\sum(q\times p)_z$守恒；单个电子碰撞中可交换角动量，不能声称每个$L_i$均守。局域碰撞总机械动量守恒、且$A(q)$相同，故也保持总canonical动量矩和$L_z$。一般热平衡$f_s\propto\exp[-\beta(H_s-\Omega L_z)-a_s]$，两数各固定、共同$\beta,\Omega$；若容器还平移对称可加相应总动量乘子。

固定杂质一般吸收动量与角动量，$\Omega$不再守恒；非磁杂质仍保两自旋数，磁散射可使自旋翻转，相关化学势差驰豫。自旋数差给额外连续/扩散模式，角动量守恒通常由动量方程乘位置导出，不独立增加一个流体变量；在有外部力矩或自旋--轨道耦合时需重新列守恒量。
\end{solution}
\begin{exercise}[原题6(a)--(j)]\label{r:333-2-6}
二维Lorentz气，固定硬圆散射体密度$n_0$、半径$a$。(a)角$\theta\in[0,2\pi]$微分截面$a\sin(\theta/2)\dd\theta/2$与总截面；(b)仅杂质Boltzmann；(c)任意$\chi(|p|)$的矩方程；(d)质量连续；(e)能量方程用$\epsilon=\rho\avg{c^2}/2,c=p/m-u,h=\rho\avg{cc^2}/2,P_{ij}=\rho\avg{c_ic_j}$；(f)$f_0=n\e^{-p^2/(2m\kb T)}/(2\pi m\kb T)$的零阶矩；(g)单碰撞时间一阶$f$；(h)质量流；(i)$\chi_{ij}=\partial u_i/\partial F_j$；(j)压力、热流、一阶流体式。
\end{exercise}
\begin{solution}
冲击参数有符号$b=a\cos(\theta/2)$，全角$|\dd b|=(a/2)\sin(\theta/2)\dd\theta$，总$2a$。若只用无符号偏转$0\le\theta\le\pi$，两侧合计为$a\sin(\theta/2)\dd\theta$，不能把两个约定再双倍。

$Lf=n_0|p|/m\int\dd\theta\,(\dd\sigma/\dd\theta)[f(p')-f(p)]$。因$|p'|=|p|$，任意径向$\chi$的碰撞积分0，分部积分给$\partial_t(n\avg\chi)+\nabla\cdot[n\avg{(p/m)\chi}]=nF\cdot\avg{\nabla_p\chi}$。$\chi=1$得$\partial_t\rho+\nabla\cdot\rho u=0$；能量得
\begin{equation}
\partial_t(\epsilon+\rho u^2/2)+\nabla\cdot[u(\epsilon+\rho u^2/2)+h+P\cdot u]=nF\cdot u.
\end{equation}
纯Lorentz气还保每个$|p|$分布，故需弱粒子碰撞才可用仅数、能的局部Maxwell闭合。零阶$u=h=0$、$P_{ij}=n\kb T\delta_{ij}$、$\epsilon=n\kb T$，各局域数能暂时无驰豫。

令$C=-(f-f_0)/\tau$，准静态一次梯度$f=f_0[1-\tau(p/m)\cdot(\nabla\log\rho-\nabla\log T+p^2\nabla T/(2m\kb T^2)-F/(\kb T))]$。Gaussian矩$\avg{p_ip_j}=s\delta_{ij}$、$\avg{p_ip_jp^2}=4s^2\delta_{ij}$、$\avg{p_ip_jp^4}=24s^3\delta_{ij}$，$s=m\kb T$。所以
\begin{equation}
\rho u=n\tau[F-\kb T\nabla\log(\rho T)],\quad \chi_{ij}=\tau\delta_{ij}/m,\quad P_{ij}=n\kb T\delta_{ij}+O(\nabla^2).
\end{equation}
实验室能流$Q=2n\kb Tu-2n\tau\kb^2T\nabla T/m$。转到题定义的\textbf{peculiar}热流要减$2n\kb Tu$，故$h=-2n\tau\kb^2T\nabla T/m$。把这些代上述质量、能量式就是一阶流体闭合，无独立保动量Euler方程。

官方解题6(j)把二维Wick六阶系数写22并给出有量纲问题的热流；独立以$\avg{p^6}=2\cdot4\cdot6\,s^3=48s^3$、分量各半验证应为24，且必须减对流焓流。本题据矩积分修复，完整有限数值核对随QA给出，不能将官方误植再当正确结果。
\end{solution}
\begin{exercise}[原题7(a)--(e)]\label{r:333-2-7}
三维经典气夹于$y=0,w$两板，温$T_1,T_2$，无宏观速，$f_0=n(y)\e^{-p^2/(2m\kb T(y))}/[2\pi m\kb T(y)]^{3/2}$。(a)$n,T$必要关系；(b)证明$\avg{p^2}=3s,\avg{p^4}=15s^2,\avg{p_y^2p^4}=35s^3$，可用$\avg{p^6}=105s^3$；(c)一碰撞时长$\tau_K$的$f_1$；(d)热导率；(e)稳态$T(y)$。
\end{exercise}
\begin{solution}
无力动量平衡要求压强$n\kb T$常数，因此$n'/n=-T'/T$。独立Gaussian三分量的Wick矩或$\chi^2_3$径向积分给3、15、105，旋转均分六阶为35。流$f_1=f_0-\tau_K(p_y/m)\partial_yf_0=f_0[1-\tau_K(p_y/m)(T'/T)(p^2/(2m\kb T)-5/2)]$，该5/2保证质量流0。

$h_y=\int(p_y/m)(p^2/2m)f_1\dd^3p$，五阶项奇为0，六阶和四阶差给$h_y=-5n\tau_K\kb^2T\,T'/(2m)$，热导$\kappa=5n\tau_K\kb^2T/(2m)$。稳态热流常量，所以$\int_{T_1}^{T(y)}\kappa(T)\dd T=(y/w)\int_{T_1}^{T_2}\kappa(T)\dd T$。若题模型$\tau_K$为常数且$nT$常数，$\kappa$恒定，温度线性；若硬球实际$\tau_K\propto1/(n\sqrt T)$，则$\kappa\propto\sqrt T$，$T(y)^{3/2}$线性。原题未规定真实碰撞截面的温变时，不能不加条件宣称唯一线性或非线性剖面。
\end{solution}
